%!TEX root = main.tex

\section{First Section}

\begin{frame}{}
    \frametitle[Short Frame Title]{Full Frame Title}
    \framesubtitle{subtitle}

    Lorem ipsum dolor sit amet, consectetur adipisicing elit, sed do eiusmod
    tempor incididunt ut labore et dolore magna aliqua. Ut enim ad minim veniam,
    quis nostrud exercitation ullamco laboris nisi ut aliquip ex ea commodo
    consequat.\footnote{Alpha}\footnote{Beta}\ufootnote{Gamma}

    \begin{equation}
        E = mc^{2}
    \end{equation}

    Duis aute irure dolor in \alert{reprehenderit} in voluptate velit esse
    cillum dolore eu fugiat nulla pariatur. Excepteur sint occaecat cupidatat non
    proident, sunt in culpa qui officia deserunt mollit anim id est laborum.
\end{frame}

\begin{frame}[plain]{Second Frame}
    Lorem ipsum dolor sit amet, consectetur adipisicing elit, sed do eiusmod tempor incididunt ut labore et dolore magna aliqua.

    \begin{itemize}
        \item ABC
        \item ABC
        \begin{itemize}
            \item DEF
            \begin{itemize}
                \item XYZ
                \item XYZ
            \end{itemize}
            \item DEF
        \end{itemize}
        \item ABC
    \end{itemize}

    \pause

    Lorem ipsum dolor sit amet, consectetur adipisicing elit, sed do eiusmod tempor incididunt ut labore et dolore magna aliqua.
\end{frame}

\section{Second Section}

\begin{frame}{Third Frame}
    Lorem ipsum dolor sit amet, consectetur adipisicing elit, sed do eiusmod
    tempor incididunt ut labore et dolore magna aliqua. Ut enim ad minim veniam,
    quis nostrud exercitation ullamco laboris nisi ut aliquip ex ea commodo
    consequat. Duis aute irure dolor in reprehenderit in voluptate velit esse
    cillum dolore eu fugiat nulla pariatur. Excepteur sint occaecat cupidatat non
    proident, sunt in culpa qui officia deserunt mollit anim id est laborum.

    \begin{enumerate}
        \item ABC
        \item ABC
        \begin{enumerate}
            \item DEF
            \item DEF
            \begin{enumerate}
                \item XYZ
                \item XYZ
            \end{enumerate}
        \end{enumerate}
        \item ABC
    \end{enumerate}
\end{frame}

\begin{frame}
    Lorem ipsum dolor sit amet, consectetur adipisicing elit, sed do eiusmod tempor incididunt ut labore et dolore magna aliqua.

    \begin{block}{Theorem}
        Lorem ipsum dolor sit amet, consectetur adipisicing elit, sed do eiusmod tempor incididunt ut labore et dolore magna aliqua.
    \end{block}

    Lorem ipsum dolor sit amet, consectetur adipisicing elit, sed do eiusmod tempor incididunt ut labore et dolore magna aliqua.
\end{frame}

\section{Third Section}

\begin{frame}{Comparison of Moduli Spaces}
  \begin{columns}[T]
    \begin{column}{0.48\textwidth}
      \begin{table}
        \centering
        \begin{tabular}{lcc}
          \toprule
          \textbf{Space} & \textbf{Dimension} & \textbf{Compact?} \\
          \midrule
          $M_g$ (smooth) & $3g-3$ & No \\
          $\overline{M}_g$ (stable) & $3g-3$ & Yes \\
          $M_{g,n}$ (marked) & $3g-3+n$ & No \\
          $\overline{M}_{g,n}$ & $3g-3+n$ & Yes \\
          \bottomrule
        \end{tabular}
        \caption{Dimensions of moduli spaces}
      \end{table}
    \end{column}
    
    \begin{column}{0.48\textwidth}
      \begin{exampleblock}{Example: $g=1$ (Elliptic Curves)}
        For genus 1 curves with a marked point:
        \begin{itemize}
          \item $M_{1,1} \cong \mathcal{H}/\operatorname{SL}(2,\Z)$
          \item Dimension = 1
          \item Compactification adds nodal cubic
        \end{itemize}
      \end{exampleblock}
    \end{column}
  \end{columns}
  
  \vspace{1em}
  \begin{block}{Remark}
    The boundary $\overline{M}_g \setminus M_g$ consists of nodal curves with lower genus components.
  \end{block}
\end{frame}